%% tblnotes.tex -- documentation for the tblnotes package
%%
%% Copyright (c) 2026 Munehiro Yamamoto <munepixyz@gmail.com>
%% License: LPPL 1.3c or later
%%
%% Build: pdflatex tblnotes.tex (three times)
\documentclass{article}
\usepackage[margin=28mm]{geometry}
\usepackage{xcolor}
\usepackage{longtable,float}
\usepackage{tblnotes}
\usepackage[colorlinks, allcolors=blue!50!black]{hyperref}
\newcommand*\tnversion{v1.1 (2026/11/16)}
\newcommand*\cs[1]{\texttt{\char`\\#1}}
\newcommand*\pkg[1]{\textsf{#1}}
\newcommand*\key[1]{\texttt{#1}}
\newcommand*\meta[1]{$\langle$\textit{#1}$\rangle$}

\title{The \pkg{tblnotes} package\\
  \large Footnotes in tables as table notes or page footnotes}
\author{Munehiro Yamamoto\\
  \url{https://github.com/munepi/tblnotes}}
\date{\tnversion}

\begin{document}
\maketitle

\section{Introduction}

A \cs{footnote} inside a \texttt{tabular} gives its mark, but its text
is lost: the footnote insertion is trapped in the box of the table.
\texttt{longtable} puts the texts at the bottom of the page.  What a
book usually wants is either
\begin{itemize}
\item \emph{table notes}: the notes put right below the table, in the
  width of the table, marked a, b, \dots{} and numbered afresh for
  each table; or
\item \emph{page footnotes}: the notes put at the bottom of the page,
  numbered along with the footnotes of the body text.
\end{itemize}
The \pkg{tblnotes} package does both, with the ordinary \cs{footnote}
in the table: no special command for the mark, and no list of notes to
write by hand.  In the outermost \texttt{tabular}, \texttt{tabular*} or
\texttt{longtable}, the footnote texts are collected, together with
their marks, while the table is typeset, and are put below the table
(\key{mode=table}, the default) or at the bottom of the page
(\key{mode=page}).

\begin{table}[ht]
\centering
\caption{Table notes (\key{mode=table})}
\begin{tabular}{lll}
\hline
Engine & Output & Note \\
\hline
pdf\LaTeX & PDF & \\
Xe\LaTeX & XDV\footnote{Converted to PDF by xdvipdfmx.} & \\
Lua\LaTeX & PDF & with Lua\TeX-ja\footnote{For Japanese.} \\
up\LaTeX & DVI\footnote{Converted to PDF by dvipdfmx.} & \\
\hline
\end{tabular}
\end{table}

\begin{table}[H]
\centering
\tblnotessetup{mode=page}
\caption{Page footnotes (\key{mode=page})}
\begin{tabular}{ll}
\hline
Mode & Notes \\
\hline
\key{table} & below the table\footnote{In the width of the table.} \\
\key{page}  & at the bottom of the page\footnote{Numbered along with
  the footnotes of the body text.} \\
\hline
\end{tabular}
\end{table}

\paragraph{Supported formats/engines}
pdf\LaTeX, Xe\LaTeX, Lua\LaTeX{} (with or without Lua\TeX-ja),
p\LaTeX{} and up\LaTeX{} (with dvipdfmx).

\section{Usage}

\begin{verbatim}
\usepackage[<options>]{tblnotes}
\tblnotessetup{<options>}

\begin{tabular}{ll}
  Item\footnote{A table note.} & Value \\
\end{tabular}
\end{verbatim}

\noindent
The options are given as package options or with \cs{tblnotessetup},
which can also be used locally, e.g.\ inside a \texttt{table}
environment before the table.

\begin{description}
\item[\key{mode=table}] (default) The footnote texts are put below the
  table, in the width of the table, as the footnotes of a
  \texttt{minipage} (\cs{footnotesize}, \cs{@makefntext}).  The marks
  are made with the counter given by \key{counter} and
  \cs{the}\meta{counter}.
\item[\key{mode=page}] The footnote texts are put at the bottom of the
  page, as ordinary footnotes numbered with the \texttt{footnote}
  counter.  For a \texttt{longtable}, this is what \texttt{longtable}
  does by itself.
\item[\key{mode=off}] The tables are left as they are.
\item[\key{counter=}\meta{counter}] The counter of the table notes
  (\key{mode=table}); the default is \texttt{mpfootnote}, as in a
  \texttt{minipage}.  It must be defined by \cs{newcounter}.  For
  example:
\begin{verbatim}
\newcounter{tblnote}
\renewcommand\thetblnote{\fnsymbol{tblnote}}
\tblnotessetup{counter=tblnote}
\end{verbatim}
\item[\key{reset=true}] (default) The counter of the table notes is
  reset to zero at the beginning of each table.  With
  \key{reset=false}, the numbering goes on from table to table.
\item[\key{rule=false}] (default) No \cs{footnoterule} above the table
  notes; \key{rule=true} puts it.
\item[\key{width=}\meta{dimen}] The width of the table notes
  (\key{mode=table}).  By default (\key{width=} with an empty value),
  the width of the table is measured: the width of its box, which
  includes the \cs{tabcolsep} at both ends unless \verb|@{}| removes
  it.  The left edge of the notes is that of the table, and the table
  does not move: notes wider than the table stick out to the right.
  Give it locally before a table, e.g.\
  \verb|\tblnotessetup{width=\linewidth}|.
\item[\key{skip=}\meta{glue}] The space between the table and the table
  notes (it sets \cs{skip}\cs{@mpfootins} for them).  By default the
  space of the footnotes of a \texttt{minipage} is used.
\end{description}

\subsection{What is collected}

\begin{itemize}
\item \cs{footnote}, \cs{footnote}\texttt{[\meta{number}]},
  \cs{footnotemark} and \cs{footnotetext} in the table, in any column
  type (including the \texttt{p} columns of \texttt{longtable}).  With
  \key{mode=table}, \cs{footnotemark} uses the counter of the table
  notes.
\item The footnotes of a table inside the table: the outermost table
  collects them.
\item The footnotes in the head of a \texttt{longtable}: the text is
  collected once, and the mark is repeated on every page.
\end{itemize}
Inside a \texttt{minipage}, a \pkg{tcolorbox} and the like, where the
footnotes are already those of the box, the tables are left as they
are.  Inside a \pkg{tcolorbox} with the key \key{page footnote} of the
\pkg{tcbpagefootnote} package, both packages work together.

\subsection{Where the notes go}

\begin{description}
\item[\texttt{tabular}, \texttt{tabular*} (\key{mode=table})] The table
  and its notes are put into a \cs{vtop} whose baseline is that of the
  table: \texttt{[t]}, \texttt{[c]} and \texttt{[b]} and the table in a
  line of text work as before.
\item[\texttt{longtable} (\key{mode=table})] The notes are put below
  the last row, at the place given by \cs{LTleft} and \cs{LTright}, in
  the width of the table (the column widths \texttt{longtable} saves in
  the \texttt{.aux} file: run \LaTeX{} again when \texttt{longtable}
  asks for it).  A page break between the last row and the notes is
  not allowed.  With \cs{endfoot} or \cs{endlastfoot}, the notes are
  added below the last foot.
\item[\texttt{tabular}, \texttt{tabular*} (\key{mode=page})] The
  footnotes are put right after the table, in the paragraph.  For a
  table in a float, see section~\ref{sec:limitations}.
\end{description}

\subsection{hyperref}

With \pkg{hyperref}, each footnote mark links to its own footnote text:
the anchor of each mark is collected with the text.  The anchor and the
mark of the last footnote are restored after the texts are put, so that
a \cs{footnotetext} after the table for a \cs{footnotemark} in it gets
the right anchor.  The anchors are not duplicated, since the tables are
not typeset twice.

\subsection{Tagged PDF}\label{sec:tagging}

With \verb|\DocumentMetadata{tagging=on}| (Lua\LaTeX{} and
pdf\LaTeX), the footnote marks and texts are tagged by \LaTeX{} (as
\texttt{Lbl} and \texttt{FENote}), and each mark and its text refer
to each other, in the table notes as well as in the page footnotes.
\begin{itemize}
\item The table notes of a \texttt{tabular} are tagged as a
  \texttt{minipage} (\texttt{Div}) right after the table.
\item The table notes of a \texttt{longtable} are put right after the
  table (after the last foot, if any) in a paragraph of their own, at
  the place of \cs{LTpost}: inside the table, they would be tagged as a
  part of the table (or as an artifact, in the foot).  An empty box
  takes their place in the table, so that the table and its notes are
  kept on the same page, as without tagging.  Its height is that of
  the notes on the previous run, saved in the \texttt{.aux} file, as
  \texttt{longtable} saves its column widths: run \LaTeX{} again when
  \pkg{tblnotes} asks for it.
\item Inside a \texttt{minipage}, where \cs{@footnotetext} is not
  \cs{@mpfootnotetext} with tagging, the tables are left as they are
  as well.
\end{itemize}

\section{Errors and warnings}

\begin{description}
\item[\texttt{Package '<package>' cannot be used with tblnotes.}]
  The packages \pkg{threeparttable} (also loaded by
  \pkg{threeparttablex}) and \pkg{tablefootnote} do the same job by
  other means: they cannot be loaded with \pkg{tblnotes}, before or
  after it.  Write
  \cs{footnote} in the table as usual instead.
\item[\texttt{Counter '<counter>' is not defined.}] The value of
  \key{counter} is not a counter; \texttt{mpfootnote} is used.
\item[\texttt{The footnotes of a table in this '<float>' are put \dots}]
  (warning) See section~\ref{sec:limitations}.
\item[\texttt{The height of the table notes of longtable \dots}]
  (warning, tagged PDF) Run \LaTeX{} again: see
  section~\ref{sec:tagging}.
\end{description}

\section{Limitations}\label{sec:limitations}

\begin{itemize}
\item \key{mode=page}: the footnote insertions in a float are lost, so
  the footnotes of a table in a float are put right after the float
  environment, in the text.  If the float moves to another page, they
  are not on the page of the float, and the package warns.  No warning
  is issued for \texttt{[H]} of the \pkg{float} package.
\item The table notes of a \texttt{longtable} are a single block: they
  are not broken across pages.  When they are added to a last foot
  that does not fit on the last page, \texttt{longtable} puts the foot
  on a new page.
\item The notes are set in the width of the table: those of a very
  narrow table are narrow, too (use \key{width}).
\item \cs{footnotetext} outside the table for a \cs{footnotemark} in the
  table works with \key{mode=page} only.
\item Only \texttt{tabular}, \texttt{tabular*} and \texttt{longtable}
  are supported.  \texttt{tabularx}, \texttt{xltabular},
  \texttt{tabulary} (which typeset the table more than once),
  \pkg{tabularray} and \pkg{nicematrix} (which have their own notes)
  are not.
\item Vertical writing (tate) is outside the scope of this version.
\item Tagged PDF: the height of the table notes of a
  \texttt{longtable} is saved in the \texttt{.aux} file; run \LaTeX{}
  again when \pkg{tblnotes} asks for it (see
  section~\ref{sec:tagging}).
\end{itemize}

\section{Tests}

\verb|make test| compiles the test files twice with pdf\LaTeX,
Xe\LaTeX, Lua\LaTeX, p\LaTeX{} and up\LaTeX{} (+ dvipdfmx) and checks
the log.  For every footnote, \texttt{test-common.tex} records in the
\texttt{.aux} file at shipout the page and the mark of both the footnote
mark and the footnote text, where the text is set (in the table notes,
in the full \cs{columnwidth} of the page, or in a box), its width and,
with \pkg{hyperref}, the anchors; when the \texttt{.aux} file is read
back at \verb|\end{document}|, it checks that every text is typeset
exactly once, with its own mark and anchor, where it is expected
(\texttt{TN-RESULT OK}).

\begin{center}
\begin{tabular}{ll}
  \hline
  Test & What it covers\\
  \hline
  \texttt{test-basic.tex}     & \key{mode=table}: \texttt{tabular},
                                \texttt{tabular*}, nesting, counters,
                                \key{width}\\
  \texttt{test-page.tex}      & \key{mode=page}: in the text, in floats,
                                the warning\\
  \texttt{test-longtable.tex} & \texttt{longtable} over pages, head and
                                foot, both modes\\
  \texttt{test-minipage.tex}  & \texttt{minipage}, \pkg{tcolorbox},
                                \pkg{tcbpagefootnote}\\
  \texttt{test-hyperref.tex}  & the tests above with \pkg{hyperref}\\
  \texttt{test-japanese.tex}  & Japanese classes (ltjbook, jbook,
                                ujbook)\\
  \texttt{test-errors.tex}    & the error messages\\
  \texttt{test-tagging.tex}   & tagged PDF: the structure tree
                                (\texttt{show-pdf-tags})\\
  \hline
\end{tabular}
\end{center}

\section{History}

\begin{description}
\item[v1.1 (2026/11/16)] Tagged PDF (\key{tagging=on} of \cs{DocumentMetadata}):
  the table notes of a \texttt{longtable} are put after the table, so
  that each footnote mark and its text refer to each other; the tables
  in a \texttt{minipage} are left as they are (see
  section~\ref{sec:tagging}).
\item[v1.0 (2026/11/14)] First public release.
\end{description}

\section{License}

This package is licensed under the conditions of the \LaTeX{} Project
Public License, either version~1.3c of this license or any later
version.

\medskip
\noindent
Copyright \copyright{} 2026 Munehiro Yamamoto

\end{document}
